\chapter{Practical workflow}
\label{app:workflow}
\begingroup\emergencystretch=3em

This appendix explains step by step how the code works and how the steps are connected, so that somebody else can
repeat the work. The steps are the same as in \fref{fig:pipeline}. Here I add the names of the scripts and of the
files that go from one step to the next. The code is in my repository
\url{https://github.com/Zainhanif700/BuildingControlCFD}. Its \texttt{README.md} has the commands with all options, and the
file \texttt{CHANGELOG.md} lists every change with its date. The versions are named by their date (for example
\texttt{v2026-11-23}), and a script that I changed has an ``Updated'' line with the date in its first lines. With this,
nobody has to guess whether a copy of the code is the current one. The reasons for the choices in the steps are given in the
chapters: the room and the airflow simulations in \cref{ch:room}, the \cotwo{} solver and the dataset in \cref{ch:dataset},
the forecast model in \cref{ch:model}, the results in \cref{ch:results}, the computing times in \cref{ch:notused} and
the sensor model in \cref{ch:sensors}.

\section{The idea: every step reads files and writes files}
The whole work is a chain. Every script reads the files that the step before it has written, and it writes the files for
the next step. The steps are connected only through these files. Because of this, a step can be started again alone
(for example the training with other options) without repeating the steps before it, and every result can be traced back to
the file with the window settings. \fref{fig:codeflow} shows the chain. The blue boxes are programs, the green boxes are
files or folders.

\begin{figure}[h]
  \centering
  \includegraphics[width=\textwidth]{code_flow.png}
  \caption[Code flow: scripts and files.]{Code flow. Blue: programs. Green: the files and folders that one program writes and
  the next one reads. The upper two rows are the data (once for every window setting), the right column is the
  forecast model, and the lower left row is the sensor model, which is trained on the same transition files.}
  \label{fig:codeflow}
\end{figure}

\noindent The folders are in my repository: the OpenFOAM cases, the CO$_2$ solver, the forecast model and its training are in
\texttt{experiments/\allowbreak gnot\_\allowbreak openfoam\_\allowbreak rans}, the scripts that make the dataset are in
\texttt{experiments/\allowbreak co2\_\allowbreak sampling}, and the sensor model is in \texttt{experiments/\allowbreak sensors}.
The test series of the thesis (mesh study, comparison of the inputs, learning curve and others) are scripts in the
folders \texttt{runs}. Each of them says in its first lines what it does.

\section{Step 1: the window settings}
A ``setting'' is one combination of open windows and doors. The file \texttt{scenarios\_v4.txt} has one line per setting:
the name, whether it is for training or for testing, the air speed at each of the 8 windows (a positive number is air
coming in, 0 is a closed window, $-1$ is a window that lets air out) and the state of the two doors. There are 68 settings:
N01--N28 (only the windows W1, W3 and W6) and A01--A40 (all windows and both doors); 52 are for training and 16 for testing.
The settings are drawn by the small program \texttt{scenarios\_v4.py}. The choice of the settings is described in \cref{ch:room}. Everything that follows is done for the names in this
file. Two older files, \texttt{scenarios.txt} and \texttt{scenarios\_v2.txt}, are the first and the second setup; each file
says in its first line which setup it is.

\begin{lstlisting}[breaklines=true,postbreak=\space]
cd experiments/gnot_openfoam_rans
python3 scenarios_v4.py        # writes scenarios_v4.txt
\end{lstlisting}


\section{Step 2: the airflow (OpenFOAM)}
\textbf{In:} the name of one setting and the number of cores. \textbf{Out:} a folder \texttt{cases/\allowbreak <name>ws1\_\allowbreak dx0.08} with the
air velocity saved over time, and an empty file \texttt{done} when everything has finished.

The script first lets \texttt{make\_rans\_case.py} write the OpenFOAM case for this setting: the room, the $k$-$\omega$ SST
turbulence model, the cells of 8~cm, and the seated people as the source of CO$_2$. Then \texttt{run\_rans\_case.sh} makes the
mesh (\texttt{blockMesh}, the columns are cut out of the room, the windows and doors become boundary patches, the seating
zone is marked), checks it with \texttt{checkMesh}, splits it over the cores, runs the solver \texttt{pimpleFoam} and joins the
result again. The first 150~s are saved every second, because this is the real start-up of the flow. Then the script
\texttt{continue\_case.sh} continues the same run up to 600~s and saves every 10~s. The first part is marked with the file
\texttt{done150}, and the whole case with \texttt{done}. If the script is started again, it skips what is finished and
continues a stopped run where it stopped. This is the longest step: 15--90 hours on one core for 600~s (measured), 2.7 times
faster on 4 cores (see \cref{ch:notused}). The room, the mesh and the boundary conditions are described in
\cref{ch:room}. The command for one setting on 4 cores:

\begin{lstlisting}[breaklines=true,postbreak=\space]
cd experiments/co2_sampling
bash run_full_case.sh N01 4            # writes cases/N01ws1_dx0.08/ and the file done
TRACE=1 bash run_full_case.sh N01 4     # the same, with every command in the log (if it stops without a message)
\end{lstlisting}


\section{Step 3: the CO$_2$ on the GPU}
\textbf{In:} the case folder of step 2. \textbf{Out:} the file \texttt{data\_v4/<name>.npz}.

OpenFOAM gives the air velocity, and my own CO$_2$ solver (\texttt{fv\_cons.py}, written in PyTorch) computes from it how the CO$_2$ moves.
It uses the real flow for the first 150~s and then the mean flow of the time from 150 to 600~s. It computes 30~minutes and
saves the CO$_2$ in the whole room every 30~s. The program also compares its CO$_2$ with the CO$_2$ that OpenFOAM computed itself,
at 150~s and at 600~s, as a check. One setting takes 1.25--4 hours on the GPU. The file also has the window settings, because
the model needs them as input. Why the \cotwo{} is computed separately, why the flow is saved every second for the first 150~s
and why the mean flow is used afterwards, is explained in \cref{ch:dataset}.

The first 28 settings (N01--N28) were made in two stages, because I found the problem of the save interval only after the first
dataset (\cref{ch:dataset}, section ``The corrected dataset''). First OpenFOAM ran for 600~s with a save every 10~s
(\texttt{run\_batch.sh}) and the first dataset was built from it. Then OpenFOAM ran again for the first 150~s with a save every
second, and \texttt{build\_v3.py} built the corrected dataset: the real flow for 0--150~s, then the mean flow of the first
dataset. The method is the same as in \texttt{build\_v4.py}. The settings A01--A40 were run in one go with
\texttt{run\_full\_case.sh}, because the problem was already known.

The command for one setting (it builds the dataset file and, in the same call, the transition file of step 4):

\begin{lstlisting}[breaklines=true,postbreak=\space]
cd experiments/co2_sampling
python3 build_v4.py --name A01 --case cases/A01ws1_dx0.08 --out data_v4 \
    --scen ../gnot_openfoam_rans/scenarios_v4.txt --trans-out data_transitions_v4
\end{lstlisting}


\section{Step 4: the transition files}
\textbf{In:} the file of step 3 and the file of a second setting (the ``partner'', drawn from all 68 settings with a fixed
random seed). \textbf{Out:} the file \texttt{data\_transitions\_v4/<name>.npz}.

The transition runs are explained in \cref{ch:model}. As in the paper, a transition run starts from the CO$_2$ of another setting: the window setting is changed, the people are
switched off, and the new flow carries the CO$_2$ away. The program takes the 30-minute CO$_2$ of the partner as the start and lets
my CO$_2$ solver run on the flow of this setting. The result is the kind of data the model has to learn: how the CO$_2$ changes
after the windows have changed. The CO$_2$ equation is linear, so the number of people can be chosen freely when the model
is trained. This step needs only the dataset files, no OpenFOAM case, and takes minutes. The script \texttt{collect\_v4.sh} does
steps 3 and 4 automatically for all finished cases, which is useful when the cases ran on several computers. The transition file alone (the dataset files already exist):

\begin{lstlisting}[breaklines=true,postbreak=\space]
python3 build_v4.py --name A01 --out data_v4 --scen ../gnot_openfoam_rans/scenarios_v4.txt \
    --trans-out data_transitions_v4 --only-transition
\end{lstlisting}


\section{Step 5: training the model}
\textbf{In:} the transition files, the file \texttt{scenarios\_v4.txt} (it says which setting is for training and which for
testing). \textbf{Out:} the folder \texttt{checkpoints/<tag>/} with the five trained networks \texttt{member0.pth} to \texttt{member4.pth}.

The model and its input are described in \cref{ch:model}. The program cuts the transition files into samples: the CO$_2$ maps of 6~minutes (12 maps) go in, the maps of the next
3~minutes (6 maps) come out. The final model is trained with the options \texttt{-{}-transitions -{}-open-flags -{}-pad-prob 0.9
-{}-tau -{}-iters 40000 -{}-z 1.6 -{}-c-offset 400}. They mean: use the transition files, give the model the open windows and
doors as extra input, let 90~\% of the samples have a window change inside the 6~minutes (and tell the model the time since
this change), train 40,000 steps, use the maps at 1.6~m height and add the 400~ppm of the outdoor air. A network needs about
38~minutes on an RTX A2000, and the five networks are averaged. The program prints the error on the 16 test settings at the
end. It does not overwrite old results: a \texttt{<tag>} that is already used stops the program. The command of the final model:

\begin{lstlisting}[breaklines=true,postbreak=\space]
cd experiments/gnot_openfoam_rans
python3 train_forecast.py --data-dir data_transitions_v4 --transitions \
    --scenarios scenarios_v4.txt --tag v4_tr_tau09_long --z 1.6 --c-offset 400 \
    --open-flags --pad-prob 0.9 --tau --iters 40000
\end{lstlisting}


\section{Step 6: testing the model}
\textbf{In:} the folder of the networks and the transition files. \textbf{Out:} printed errors (and pictures).

The results are in \cref{ch:results} (the window change inside the history in \cref{sec:winchange}).
\texttt{eval\_history\_change.py} builds test samples in which the window setting changes inside the 6-minute history, because
in a real room a window can be opened at any time. It prints the error of the model and of the simple guess that the CO$_2$ stays as
it is now. \texttt{oracle\_mean\_bound.py} computes in about one minute what a simple mixing-box rule without any model can
reach, as a reference. The small programs \texttt{test\_fv\_cons.py}, \texttt{test\_forecast.py} and \texttt{test\_sensors.py}
test the solver and the models with small examples and run on a CPU in a few minutes. I run them first when I change the code.

\begin{lstlisting}[breaklines=true,postbreak=\space]
python3 eval_history_change.py --ckpt checkpoints/v4_tr_tau09_long --data-dir data_transitions_v4 \
    --scenarios scenarios_v4.txt
python3 oracle_mean_bound.py --data-dir data_transitions_v4 --scenarios scenarios_v4.txt
python3 test_fv_cons.py && python3 test_forecast.py     # small tests, CPU, a few minutes
\end{lstlisting}


\section{Step 7: the sensor model}
The sensor model is described in \cref{ch:sensors}. It forecasts the whole room from the 8 sensor stations only. It is trained on the \emph{same} transition files as
the final model, and the simulated CO$_2$ at the station positions plays the role of the sensor values. The files have
the following tasks.
\begin{itemize}
  \item \texttt{stations.py}: the 8 stations (6 windows, 2 doors) with their assumed positions.
  \item \texttt{synthetic.py}: makes a sensor file in the format of the sensor server from one simulation file, with noise.
  \item \texttt{readers.py}: reads such a file (server export, older CSV or database).
  \item \texttt{mapper.py}: turns the records into the model input: 12 CO$_2$ values per station (6~minutes, every 30~s), the 8 window
        speeds and the 2 door states.
  \item \texttt{train\_sensor.py}: trains the sensor model (5 networks, noise of 10~ppm added to the station values) and
        writes the networks to \texttt{sensors/checkpoints/<tag>/}.
  \item \texttt{diagnose\_sensor.py}: prints the errors, also for the simple guesses (``stays as it is now'', ``the whole room at the
        mean of the stations''), near and far from a station.
  \item \texttt{sensor\_forecast.py}: the whole chain in one program: sensor file $\rightarrow$ reader $\rightarrow$ mapper $\rightarrow$ sensor model
        $\rightarrow$ forecast of the whole room $\rightarrow$ the CO$_2$ at a point that I choose. It draws one overview picture.
        When the real stations deliver data, only the first step (the reader of the file) changes.
  \item \texttt{warning\_rules.py} and \texttt{lecture\_demo.py}: the warning at 700 and 800~ppm, and a synthetic lecture through the whole chain.
\end{itemize}

\begin{lstlisting}[breaklines=true,postbreak=\space]
cd experiments/sensors
python3 train_sensor.py --data-dir ../gnot_openfoam_rans/data_transitions_v4 --transitions \
    --scenarios ../gnot_openfoam_rans/scenarios_v4.txt --tag my_run --noise 10
python3 diagnose_sensor.py --data-dir ../gnot_openfoam_rans/data_transitions_v4 --transitions \
    --scenarios ../gnot_openfoam_rans/scenarios_v4.txt --tag my_run
python3 sensor_forecast.py --tag st_lecture --sim ../gnot_openfoam_rans/data_v4/N24.npz --xy 7.8 4.6
\end{lstlisting}


\section{Using the trained model}
To forecast, the model needs the CO$_2$ maps of the last 6~minutes, the window speeds, the open windows and doors and the number
of people (final model), or the station values, speeds and doors (sensor model). It returns the CO$_2$ of the next 3~minutes with an
uncertainty. In the repository, the small folder \texttt{demo} contains the trained final model and one test setting (N14). The command
\texttt{bash run\_demo.sh forecast}, in \texttt{experiments/\allowbreak gnot\_\allowbreak openfoam\_\allowbreak rans},
runs the forecast, prints its error next to the error of the simple guess, and saves a picture. The command
\texttt{bash run\_demo.sh openfoam} runs step 2 for one setting for only 20~s (about 10~minutes on 4 cores), as a test of the OpenFOAM
installation. The two parts do not depend on each other.

\begin{lstlisting}[breaklines=true,postbreak=\space]
cd experiments/gnot_openfoam_rans
bash setup_env.sh              # once: installs the packages, prints python OK and OpenFOAM OK
bash run_demo.sh forecast        # forecast with the final model, error against "stays as it is now", picture
bash run_demo.sh openfoam        # OpenFOAM, N01, 20 s, 4 cores
\end{lstlisting}

\section{Overview of the files}
\begin{table}[h]
  \centering
  \caption[Files that go from step to step.]{Files that go from step to step. The times are the ones I measured, for one setting.}
  \label{tab:flowfiles}
  \footnotesize
  \begin{tabular}{p{0.4cm}p{2.9cm}p{3.2cm}p{3.2cm}p{2.4cm}}
    \toprule
    & Script & Reads & Writes & Time\\
    \midrule
    1 & \texttt{scenarios\_\allowbreak v4.py} & -- & \texttt{scenarios\_\allowbreak v4.txt} & --\\
    2 & \texttt{run\_\allowbreak full\_\allowbreak case.sh} & name of the setting & \texttt{cases/\allowbreak <name>ws1\_\allowbreak dx0.08} & 15--90~h, 1 core\\
    3 & \texttt{build\_v4.py} & the case folder & \texttt{data\_v4/\allowbreak <name>.npz} & 1.25--4~h, GPU\\
    4 & \texttt{build\_v4.py} \texttt{-{}-trans-out} & files of step 3 & \texttt{data\_\allowbreak transitions\_\allowbreak v4/\allowbreak <name>.npz} & minutes\\
    5 & \texttt{train\_\allowbreak forecast.py} & transition files, settings & \texttt{checkpoints/\allowbreak <tag>/\allowbreak member0-4.pth} & 38~min per network\\
    6 & \texttt{eval\_\allowbreak history\_\allowbreak change.py} & networks, transition files & printed errors & minutes\\
    7 & \texttt{train\_\allowbreak sensor.py}, \texttt{diagnose\_\allowbreak sensor.py} & transition files, settings & \texttt{sensors/\allowbreak checkpoints/\allowbreak <tag>/}, printed errors & see \cref{ch:sensors}\\
    \bottomrule
  \end{tabular}
\end{table}

\section{What is needed}
Step 2 needs OpenFOAM v2412 from openfoam.com. The versions from openfoam.org (12, 13, 14) are a different program and do not
work with my cases. The other steps need Python with PyTorch (the packages are in \texttt{requirements.txt}; \texttt{setup\_env.sh}
installs them and checks that Python and OpenFOAM are found) and, for the CO$_2$ solver and the training, a GPU. For the
OpenFOAM part there is also a Docker image (the file \texttt{Dockerfile}). The data files and the trained networks are large and are not
in the repository, except the small demo.

\section{Where I tested it}
All steps run on our lab computer (Linux). The demo was tested there in a clean copy of the repository. The OpenFOAM part also ran on a
Mac and, in the Docker image, on a DGX Spark (the 20~s test of all 28 settings of the second setup).

\section{If something goes wrong}
\begin{itemize}
  \item If an OpenFOAM script stops without a message, I start it again with \texttt{TRACE=1} in front of the command. The log
        then shows every command with the time and the line.
  \item A finished setting gets a file \texttt{done} and is skipped when the script is started again. A stopped run continues
        where it stopped.
  \item One finished case needs about 2--4~GB. For many settings it is best to run one setting per core.
\end{itemize}
\endgroup
